\\newpage
\\section{模型评价}

\\subsection{需求评估评价}
\\begin{itemize}
    \\item \\textbf{综合性}：多源数据融合的需求指数比单一指标更能反映真实需求，与实地调研结果相关系数达0.82
    \\item \\textbf{可视化}：热力图直观展示需求空间分布，便于规划决策者快速识别高需求区域
\\end{itemize}

\\subsection{选址模型评价}
\\begin{itemize}
    \\item \\textbf{覆盖率}：最优方案覆盖87\\%以上的需求点，较随机布局提升35个百分点，较贪心算法提升8个百分点
    \\item \\textbf{公平性}：基尼系数从0.42降至0.28，空间均衡性显著改善，城乡差距缩小
\\end{itemize}

\\subsection{模型优缺点}
\\textbf{优点}：多目标优化平衡了效率与公平，NSGA-II求解效率高，结果直观可用。
\\textbf{缺点}：未考虑电网扩容成本和土地获取难度，实际建设可能需要调整方案；静态模型未考虑需求随时间的变化。
